\documentclass[12pt,a4paper]{report}

\usepackage[utf8]{inputenc}
\usepackage[T1]{fontenc}
\usepackage[spanish,es-tabla]{babel}
\usepackage{geometry}
\usepackage{graphicx}
\usepackage{booktabs}
\usepackage{longtable}
\usepackage{tabularx}
\usepackage{amsmath}
\usepackage{hyperref}
\usepackage{enumitem}
\usepackage{xcolor}

\geometry{
  left=3cm,
  right=2.5cm,
  top=2.5cm,
  bottom=2.5cm
}

\hypersetup{
  colorlinks=true,
  linkcolor=blue,
  citecolor=blue,
  urlcolor=blue
}
\Urlmuskip=0mu plus 1mu\relax

\setlist[itemize]{leftmargin=*}
\setlist[enumerate]{leftmargin=*}
\newcolumntype{Y}{>{\raggedright\arraybackslash}X}

\title{
  Evaluación del comportamiento de agentes con memoria RAG\\
  en el análisis \textit{run-to-failure} de maquinaria rotativa
}
\author{Autor: pendiente}
\date{Curso académico: pendiente}

\begin{document}

\pagenumbering{gobble}
\maketitle
\cleardoublepage{}
\pagenumbering{roman}

\chapter*{Resumen}
Este Trabajo Fin de Máster estudia el comportamiento de agentes basados en
modelos de lenguaje locales durante el análisis de señales industriales de
vibración. El escenario principal es la monitorización \textit{run-to-failure}
de rodamientos: los agentes deben formular hipótesis, seleccionar
configuraciones, interpretar resultados y revisar sus decisiones a partir de
evidencia temporal. Para que ese razonamiento pueda evaluarse, la propuesta
combina un supervisor jerárquico, decisiones validadas mediante contratos
estrictos y ejecutores deterministas que realizan las transformaciones sobre
los datos.

La aportación diferencial es un ciclo de memoria RAG supervisada. Cada decisión
puede producir un episodio trazable y un candidato de memoria; antes de volver a
utilizarse, los recuerdos se filtran por agente, dataset, rol, veredicto y
compatibilidad metodológica. El sistema registra qué recuerdos se recuperan,
cuáles llegan al agente y cómo este los cita, adapta o rechaza. El índice
vectorial es intercambiable y reconstruible; los contratos y artefactos
persistidos constituyen la fuente auditable. El almacenamiento JSON exacto
actúa como referencia reproducible y Qdrant queda como implementación opcional.

La implementación se ha validado como sistema mediante CWRU y pilotos
sintéticos, y se ha ejecutado un primer baseline causal sobre las 984
adquisiciones oficiales de NASA IMS Set~2. El protocolo fija, antes del
modelado, 197 snapshots de baseline, 98 de calibración y 689 de monitorización;
calibra el umbral sin etiquetas futuras y evalúa persistencia después de agregar
las ventanas de cada adquisición. Esta ejecución aporta evidencia descriptiva
de una trayectoria real, pero no valida un onset físico ni permite publicar
métricas binarias sin \textit{ground truth}. Un piloto pareado con PCA e
Isolation Forest recuperó un único recuerdo incompatible, que no fue citado ni
utilizado; las métricas permanecieron invariantes. Este resultado negativo
revela que la disponibilidad del motor no compensa un corpus irrelevante y no
demuestra un beneficio del RAG. El contraste definitivo requiere réplicas con
un corpus pertinente congelado y una comparación controlada de modelos de
lenguaje. La evaluación determinista previa no autorizó nuevas generaciones al
no disponer aún de trazas experimentales completas, cobertura multi-dataset y
resultados de las perturbaciones definidas.

\noindent\textbf{Palabras clave:} sistemas multiagente, modelos de lenguaje,
memoria RAG, mantenimiento predictivo, run-to-failure, detección de anomalías,
señales de vibración.

\chapter*{Abstract}
This Master's Thesis studies the behaviour of local language-model agents when
analysing industrial vibration signals. Its main scenario is bearing
run-to-failure monitoring: agents formulate hypotheses, select configurations,
interpret results and revise their decisions using temporal evidence. A
hierarchical supervisor coordinates the workflow, strict contracts validate
agent decisions, and deterministic executors perform all data transformations.

The distinctive contribution is a supervised RAG memory cycle. Decisions may
produce traceable episodes and memory candidates; before reuse, records are
filtered by agent, dataset, role, verdict and methodological compatibility. The
system records which memories are retrieved, which ones reach the agent, and
how the agent cites, adapts or rejects them. The vector index is interchangeable
and rebuildable; persisted contracts and artefacts remain the auditable source
of evidence. Exact JSON search is the reproducible baseline, while Qdrant is an
optional implementation.

The implementation has been validated as a system with CWRU and synthetic
pilots, and a first causal baseline has been run on the 984 official NASA IMS
Set~2 acquisitions. Before modelling, the protocol fixes 197 baseline, 98
calibration and 689 monitoring snapshots; it calibrates the threshold without
future labels and evaluates persistence only after aggregating each
acquisition's windows. This run provides descriptive evidence from a real
trajectory, but it neither validates a physical onset nor supports binary
metrics without ground truth. In a paired PCA and Isolation Forest pilot, the
retriever returned one incompatible record; agents neither cited nor used it,
and the metrics remained unchanged. This negative result shows that an
available retrieval engine cannot compensate for an irrelevant corpus and
does not demonstrate a RAG benefit. Definitive evidence requires replicated
ablations with a relevant frozen corpus and a controlled comparison of language
models. The deterministic pre-generation assessment did not authorise new
generations because complete experimental traces, multi-dataset coverage and
perturbation results were not yet available.

\noindent\textbf{Keywords:} multi-agent systems, language models, RAG memory,
predictive maintenance, run-to-failure, anomaly detection, vibration signals.

\tableofcontents
\listoffigures
\listoftables

\cleardoublepage{}
\pagenumbering{arabic}

\chapter{Introducción}

\section{Contexto y motivación}

El mantenimiento predictivo busca reconocer cambios de comportamiento antes de
que una avería obligue a detener un activo. En maquinaria rotativa, las señales
de vibración permiten observar fenómenos como impactos, desalineación,
desbalanceo o degradación de rodamientos. Sin embargo, transformar esas señales
en una decisión útil no consiste únicamente en entrenar un detector. También es
necesario justificar cómo se prepararon los datos, por qué se eligió un modelo,
qué evidencia respalda una alerta y qué limitaciones tiene el resultado.

Los modelos de lenguaje ofrecen una forma flexible de interpretar esa evidencia
y coordinar tareas especializadas. Su uso plantea, a la vez, una dificultad
metodológica: si el agente responde en texto libre o ejecuta directamente las
transformaciones, resulta difícil separar razonamiento, cálculo y error. Este
trabajo adopta una separación explícita. Los agentes proponen decisiones
estructuradas; los ejecutores deterministas procesan las señales; y todos los
pasos relevantes dejan artefactos que pueden revisarse después.

El interés principal no es, por tanto, obtener una única métrica de detección,
sino observar si agentes locales relativamente pequeños son capaces de
formular hipótesis razonables, utilizar herramientas, revisar resultados y
aprender de experiencias anteriores sin perder reproducibilidad.

\section{Planteamiento del problema}

El problema se aborda desde dos escenarios complementarios. CWRU Bearing Data
Center proporciona un benchmark etiquetado y acotado que permite comprobar el
pipeline y mantener una referencia de regresión~\cite{cwru_bearing_data_center}.
El escenario de investigación principal es NASA IMS Bearing, donde una secuencia
de adquisiciones describe la evolución del activo hasta el final del ensayo
~\cite{nasa_pcoe_repository}. En este segundo caso, el objetivo no debe reducirse
a clasificar ventanas aisladas: interesa estudiar la tendencia del indicador de
salud, la persistencia de las alertas, su tasa antes de monitorización y su
posición retrospectiva respecto al final registrado.

Sobre ambos escenarios se incorpora una memoria RAG\@. Un recuerdo puede ayudar a
repetir una decisión útil, advertir sobre un fallo anterior o señalar que una
estrategia no es comparable con el caso actual. El riesgo es que una memoria
irrelevante, incorrecta o autorreferencial condicione decisiones posteriores.
Por ello, recuperar el fragmento más próximo no es suficiente: el ciclo necesita
gobierno, filtrado, citas verificables y evaluación de su efecto.

\section{Preguntas de investigación}

El trabajo se organiza alrededor de las siguientes preguntas:

\begin{enumerate}[label=\textbf{PI\arabic*:}]
  \item ¿Pueden agentes basados en modelos de lenguaje locales formular y
        justificar decisiones técnicas útiles dentro de contratos estrictos,
        delegando el cálculo en ejecutores reproducibles?
  \item ¿Cómo cambia su razonamiento cuando reciben memoria de decisiones
        anteriores, y pueden citar, adaptar o rechazar esa memoria de forma
        trazable?
  \item ¿Qué mecanismos de filtrado y promoción reducen la reutilización de
        recuerdos incompatibles, no validados o perjudiciales?
  \item ¿Cómo se comporta el sistema en un problema \textit{run-to-failure},
        donde la unidad longitudinal es la secuencia de snapshots y las
        métricas operacionales difieren de una clasificación binaria?
  \item Una vez fijado el protocolo experimental, ¿qué diferencias aparecen
        entre Qwen 3.5 y Qwen 3 bajo las mismas entradas, herramientas,
        contratos y condiciones de memoria?
\end{enumerate}

La última pregunta se plantea como comparación posterior. No se deben atribuir
diferencias al modelo mientras cambien simultáneamente el dataset, la política
temporal o el corpus de memoria.

\section{Objetivos}

\subsection{Objetivo general}

Diseñar y evaluar una plataforma híbrida y auditable para monitorización
industrial \textit{run-to-failure}. La propuesta combina un observador
determinista frecuente con supervisión multiagente periódica o dirigida por
eventos, de modo que los agentes puedan diseñar y auditar políticas de
detección antes de un eventual despliegue industrial. El estudio presta
especial atención a sus hipótesis y al efecto de una memoria RAG gobernada
sobre las decisiones.

\subsection{Objetivos específicos}

\begin{enumerate}
  \item Coordinar agentes especializados mediante un supervisor jerárquico y
        representar sus decisiones con esquemas estructurados.
  \item Separar el razonamiento generativo de los ejecutores que transforman
        datos, entrenan modelos y calculan métricas.
  \item Mantener trazabilidad de entradas, configuraciones, decisiones,
        resultados, revisiones y artefactos.
  \item Implementar un ciclo de memoria que distinga captura, validación,
        promoción, recuperación, uso y auditoría.
  \item Evaluar la relevancia de la memoria y la fidelidad de las citas, además
        de su posible efecto sobre configuraciones y resultados.
  \item Definir una metodología temporal para NASA IMS que separe ventanas
        intra-snapshot de observaciones longitudinales por snapshot.
  \item Implementar un replay histórico causal que permita ensayar la
        separación entre monitorización determinista y revisión multiagente
        sin presentar la simulación como adquisición industrial en tiempo real.
  \item Comparar baselines interpretables y modelos de anomalías bajo un mismo
        protocolo temporal.
  \item Preparar una comparación controlada entre Qwen 3.5 y Qwen 3.
\end{enumerate}

\section{Aportaciones del trabajo}

Las aportaciones del trabajo se agrupan en cuatro bloques. En primer lugar, una
arquitectura híbrida que mantiene el cálculo frecuente en ejecutores
deterministas y hace observable la revisión periódica de los agentes sin
permitirles ejecutar código arbitrario. En segundo lugar, un ciclo RAG orientado
a episodios de decisión y no solo a documentos: conserva alternativas,
evidencia, resultado y condiciones de reutilización. En tercer lugar, un
protocolo de evaluación que separa calidad de recuperación, fidelidad de uso,
cambio de comportamiento y resultado operacional. Finalmente, una aplicación
del enfoque al seguimiento \textit{run-to-failure} de rodamientos mediante
replay histórico causal.

\section{Alcance y limitaciones}

El sistema se ejecuta localmente y emplea modelos disponibles mediante Ollama.
Los agentes no modifican directamente los datos ni generan código para
ejecutarlo. La arquitectura contempla distintos datasets y backends de memoria,
pero el estudio no pretende resolver despliegue multiusuario, adquisición
continua de planta, planificación distribuida o seguridad productiva. La vista
de monitorización reproduce una trayectoria histórica con un reloj acelerado;
su semejanza visual con una consola en vivo no acredita latencia, disponibilidad
ni integración con sensores industriales.

La evidencia actual debe interpretarse con cautela. CWRU valida el flujo
etiquetado y los pilotos sintéticos sirven como pruebas de integración. Sobre
NASA IMS Set~2 oficial se dispone ya de un baseline PCA causal que recorre sus
984 snapshots con partición 20/10/70, calibración separada y evaluación
longitudinal agregada por adquisición. La vista de decisión oculta a limpiador,
estructurador y modelador los campos retrospectivos. Esta única trayectoria
permite estudiar el mecanismo y describir la evolución del score, pero no
aporta etiquetas continuas de degradación, un onset físico oficial ni evidencia
de generalización. Aún es necesario congelar el corpus de memoria y ejecutar
ablaciones agénticas con réplicas.

\section{Estructura del documento}

El capítulo 2 revisa mantenimiento predictivo, agentes y memoria RAG\@. El
capítulo 3 presenta la arquitectura y los contratos. El capítulo 4 define la
metodología de evaluación agéntica, de memoria y \textit{run-to-failure}. El
capítulo 5 resume las decisiones de implementación relevantes, sin reproducir
la estructura interna del repositorio. Los capítulos 6 y 7 describen el diseño
experimental y los resultados, separando pilotos sintéticos de evidencia real.
El capítulo 8 recoge conclusiones, limitaciones y trabajo futuro.

\chapter{Estado del arte}

\section{Mantenimiento predictivo y detección de anomalías}

Esta sección desarrollará el contexto de mantenimiento predictivo, monitorización
de condición y detección de anomalías en maquinaria rotativa.

Aspectos a cubrir:

\begin{itemize}
  \item diferencias entre mantenimiento correctivo, preventivo y predictivo;
  \item papel de las señales de vibración en rodamientos y motores;
  \item fallos comunes: pista interior, pista exterior, bola y desbalanceo;
  \item diferencias entre detección de anomalías, diagnóstico y pronóstico.
\end{itemize}

\section{Modelos de detección de anomalías}

El MVP comparará primero modelos clásicos y reproducibles:

\begin{itemize}
  \item Isolation Forest;
  \item One-Class SVM;
  \item Local Outlier Factor;
  \item PCA con error de reconstrucción.
\end{itemize}

Los autoencoders densos y modelos recurrentes se considerarán en fases
posteriores, cuando el pipeline de tensores esté validado.

En la extension run-to-failure, la lectura de estos modelos cambia. Un detector
no se evalua solo por separar ventanas normales y anomalias, sino por producir
una trayectoria de score que anticipe degradacion, mantenga falsas alarmas
nominales controladas y permita construir un indicador de salud interpretable.
Por ello, PCA con error de reconstruccion, Isolation Forest, One-Class SVM y
autoencoders densos se tratan como familias comparables dentro de un mismo
protocolo temporal, no como modelos aislados optimizados por F1.

\section{Run-to-failure, indicadores de salud y modelos de reconstruccion}

Los datasets run-to-failure representan una trayectoria historica de un activo
hasta fallo. En este tipo de problema conviene distinguir cuatro conceptos:
anomalia puntual, degradacion sostenida, indicador de salud y pronostico de
vida util restante. Una alerta aislada puede deberse a ruido o transitorios; la
degradacion defendible requiere persistencia temporal, tendencia del score y
coherencia con el tramo nominal.

Los Health Indicators son habituales en mantenimiento predictivo porque
condensan la evolucion del activo en una curva de salud. En un sistema
experimental local, este indicador debe declararse como evidencia auxiliar y no
como prediccion industrial cerrada. Su calidad puede evaluarse mediante caida
entre tramo inicial y final, monotonicidad, robustez frente al ruido, volatilidad
nominal y fuerza de tendencia durante la degradacion.

Los autoencoders densos aportan una extension natural para este escenario,
porque aprenden a reconstruir normalidad y convierten el error de
reconstruccion en un score de anomalia. Sin embargo, su uso exige suficientes
ventanas nominales de entrenamiento, control de fuga temporal, arquitectura
cerrada, evaluacion comparable y guardarrailes que impidan presentar un modelo
mas complejo como mejora automatica. En este TFM se incorporan como familia
PyTorch reproducible sobre features tabulares, dejando los autoencoders
recurrentes y el RUL completo como lineas futuras.

\section{Sistemas multiagente aplicados a pipelines de datos}

Esta sección describirá el uso de arquitecturas multiagente para dividir
responsabilidades entre nodos especializados. El enfoque del proyecto separa:

\begin{itemize}
  \item decisiones de alto nivel tomadas por agentes;
  \item validación de contratos mediante esquemas estrictos;
  \item ejecución determinista de transformaciones sobre datos;
  \item auditoría de rutas, configuraciones, métricas y artefactos.
\end{itemize}

\section{Datasets públicos de referencia}

El dataset principal será CWRU Bearing Dataset
\cite{cwru_bearing_data_center,cwru_normal_baseline,cwru_12k_drive_end}. NASA
IMS Bearing Dataset se reserva como extensión para escenarios de degradación
temporal \cite{nasa_pcoe_repository,nasa_ims_bearings_zip}.

\chapter{Arquitectura propuesta}

\section{Principios de diseno}

La arquitectura se basa en cuatro principios:

\begin{enumerate}
  \item separar codigo, memoria academica y recursos;
  \item mantener los datos crudos inmutables;
  \item representar decisiones mediante contratos estructurados;
  \item ejecutar transformaciones mediante modulos Python deterministas.
\end{enumerate}

\section{Separacion entre agentes y ejecutores}

Los agentes no escriben ni ejecutan codigo arbitrario sobre los datos. Su papel
es proponer configuraciones, interpretar perfiles y decidir la siguiente fase
del flujo. Los ejecutores leen artefactos desde disco, aplican operaciones
deterministas y devuelven rutas, metricas y logs.

\section{Nodos previstos}

El grafo minimo previsto contiene:

\begin{itemize}
  \item supervisor;
  \item agente limpiador;
  \item ejecutor de limpieza;
  \item agente estructurador;
  \item ejecutor de estructuracion;
  \item agente modelador;
  \item ejecutor de modelado;
  \item evaluador;
  \item redactor;
  \item revision humana para trabajos costosos.
\end{itemize}

\section{Estado global}

El estado global de LangGraph no almacenara datasets completos. Debe contener
rutas, resumenes y artefactos:

\begin{itemize}
  \item identificador de ejecucion;
  \item contexto del proyecto;
  \item rutas de datos crudos, perfiles, datos limpios, tensores e informes;
  \item configuraciones de limpieza, estructuracion y modelado;
  \item metricas y evaluacion;
  \item errores, aprobaciones humanas y artefactos generados.
\end{itemize}

En la implementacion inicial se ha definido un modelo canonico `TFMStateModel`
validado con Pydantic y una vista `TFMState` basada en `TypedDict` para su uso
en LangGraph. Esta separacion permite validar contratos de forma estricta antes
de entregar el estado al grafo y, al mismo tiempo, mantener una representacion
ligera y serializable.

El estado inicial del caso CWRU queda fijado en la fase `dataset\_manifest`,
con el nodo siguiente `manifest\_executor`. El directorio crudo se registra como
\path{codigo/data/raw/cwru_bearing/mat}, el canal principal es `DE\_time` y la
frecuencia objetivo es 12 kHz. Las rutas de artefactos posteriores permanecen
como valores nulos hasta que los ejecutores correspondientes las generen.

\section{Contratos estructurados}

Los contratos iniciales incluyen:

\begin{itemize}
  \item `ProjectContext`, que define dominio, maquina, dataset, objetivo,
        canal principal y frecuencia objetivo;
  \item `DatasetManifestRow` y `DatasetManifest`, que formalizan el inventario
        de ficheros crudos y sus etiquetas;
  \item `DatasetProfileSummary`, que resume el perfil sin almacenar senales;
  \item `CleaningConfig`, `StructuringConfig` y `ModelingConfig`, que fijan las
        configuraciones que podran proponer los agentes;
  \item `SupervisorDecision`, `CleaningDecision`, `StructuringDecision`,
        `ModelingDecision`, `EvaluationDecision` y `ReportDecision`, que
        estructuran las respuestas de los agentes;
  \item `ExecutorResult` y resultados especializados, que normalizan las
        salidas de los ejecutores deterministas;
  \item `MetricsReport` y `EvaluationResult`, que estructuran la evaluacion;
  \item `PipelineError`, `HumanApproval` y `ArtifactRef`, que mantienen
        trazabilidad de errores, revisiones y artefactos.
\end{itemize}

Estos esquemas aplican validacion estricta y rechazan campos no declarados. Por
ejemplo, una fila normal del manifiesto debe declarar los campos de fallo como
nulos, mientras que una fila con fallo debe incluir tipo de fallo y diametro. De
forma analoga, un resultado de ejecutor con estado fallido debe incluir un error
estructurado, y una decision no terminal del supervisor debe indicar el siguiente
nodo del grafo.

\section{Trazabilidad}

Cada ejecucion debe permitir responder:

\begin{itemize}
  \item que datos se usaron;
  \item que configuracion fue validada;
  \item que ejecutor produjo cada artefacto;
  \item que metricas se obtuvieron;
  \item que limitaciones permanecen.
\end{itemize}

\chapter{Metodología}

\section{Caso de uso industrial}

El caso de uso seleccionado es la detección de anomalías en rodamientos de
motores eléctricos mediante señales de vibración. La elección permite trabajar
con datos industriales públicos, etiquetas interpretables y un flujo temporal
adecuado para ventanas, features y tensores.

\section{Pipeline de datos}

El flujo metodológico para el MVP queda definido como:

\begin{center}
\begin{minipage}{0.85\textwidth}
\begin{verbatim}
datos .mat crudos
-> manifiesto del dataset
-> extracción de señales
-> normalización de frecuencia de muestreo
-> perfilado estadístico
-> limpieza determinista
-> ventanas temporales
-> features y tensores
-> particiones train/validation/test
-> modelos de anomalías
-> evaluación
-> informe técnico
\end{verbatim}
\end{minipage}
\end{center}

\section{Generalización multi-dataset}

Tras consolidar el MVP sobre CWRU, la siguiente extensión metodológica consiste
en separar explícitamente tres niveles: el dataset crudo, el adaptador
determinista de dataset y el pipeline común. Esta separación evita que los
ejecutores dependan de nombres de fichero o estructuras particulares de un
benchmark concreto.

La generalización se plantea mediante un descriptor común de dataset, un
manifiesto común con una fila por unidad cruda procesable y un registro de
adaptadores. CWRU se mantendrá como benchmark de regresión, mientras que NASA
IMS se contempla como primer candidato para validar el sistema con series más
largas, múltiples canales y condiciones de degradación menos triviales.

La inspección local de NASA IMS muestra que el paquete se distribuye como un
contenedor anidado (`zip`, `7z` y `rar`) y que la aplicación deberá distinguir
entre detección de formato, selección de adaptador y generación efectiva de
manifiesto. En una futura interfaz, el usuario debería ver una previsualización
del descriptor del dataset, los formatos detectados, las limitaciones de
extracción y si el manifiesto puede generarse directamente o requiere preparar
una carpeta preextraída.

Como primer paso ejecutable, el manifiesto NASA IMS se restringe a carpetas ya
preextraídas. Cada snapshot temporal se representa como una fila común con
etiqueta inicial desconocida, fallo final conocido a nivel de set y metadatos
de frecuencia, canales y variante local. Esta decisión permite validar la
generalización del pipeline sin incorporar todavía dependencias de extracción
ni una política de etiquetado temporal.

Los agentes no interpretarán directamente carpetas crudas. Recibirán
descriptores, resúmenes de manifiesto y diagnósticos de calidad generados por
código determinista. De este modo pueden aumentar su capacidad de decisión
--por ejemplo, recomendando canal, ventana, familia de features o modelo-- sin
ejecutar código arbitrario ni saltarse los contratos del sistema.

\section{Manifiesto del dataset}

El primer artefacto metodológico será un manifiesto con una fila por fichero
crudo. Este manifiesto conectará los `.mat` originales con el resto del
pipeline. Las columnas previstas son:

\begin{longtable}{ll}
\toprule
Campo & Descripción \\
\midrule
file\_id & identificador del fichero original \\
dataset & nombre del dataset \\
source\_path & ruta del fichero crudo \\
label & clase `normal' o `fault' \\
fault\_type & tipo de fallo si existe \\
fault\_diameter\_inch & diámetro del defecto \\
load\_hp & carga del motor \\
rpm & régimen de giro \\
sensor\_channel & canal de sensor usado \\
source\_sample\_rate\_hz & frecuencia de muestreo original \\
target\_sample\_rate\_hz & frecuencia objetivo del MVP \\
source\_format & formato original \\
notes & observaciones \\
\bottomrule
\end{longtable}

\section{Extracción y perfilado}

El canal principal inicial será `DE\_time`, asociado al acelerómetro del lado
drive-end. El perfilado generará resúmenes estadísticos por fichero y por
etiqueta, incluyendo número de muestras, canales detectados, valores no finitos,
media, desviación típica, RMS, curtosis, factor de cresta y energía.

\section{Estructuración temporal}

La frecuencia objetivo inicial será 12 kHz. Las ventanas candidatas serán de
2048 o 4096 muestras con solapamiento del 50\%. Esta decisión permite comparar
features tabulares con modelos clásicos y conservar ventanas crudas para
autoencoders posteriores.

\section{Particiones}

Para evitar fuga de información entre ventanas, no se usará un split aleatorio
ingenuo. La estrategia inicial será:

\begin{itemize}
  \item entrenamiento con ventanas normales;
  \item validación con ventanas normales retenidas;
  \item test con ventanas normales retenidas y ventanas con fallo.
\end{itemize}

\section{Métricas}

Las métricas iniciales serán precision, recall, F1-score, ROC-AUC, PR-AUC, tasa
de falsos positivos y matriz de confusión. Desde el punto de vista industrial,
el criterio principal será mantener alta sensibilidad ante fallos sin disparar
en exceso los falsos positivos.

\chapter{Implementación}

\section{Estructura del repositorio}

La implementación se organiza en tres áreas:

\begin{itemize}
  \item `codigo/`: aplicación, tests, datos, modelos, informes y documentación
        técnica;
  \item `memoria/`: documento LaTeX del TFM;
  \item `recursos/`: PDFs, referencias y diagramas auxiliares.
\end{itemize}

\section{Estado Actual}

Hasta este punto se ha creado la estructura inicial de directorios, se han
descargado los datasets crudos, se ha definido el flujo de datos del MVP y se
han implementado el contrato inicial del estado global del grafo y los esquemas
Pydantic de entradas, decisiones y resultados. Los datos crudos se conservan en
`codigo/data/raw/`, junto con un README de trazabilidad.

\section{Estado global implementado}

El estado global se ha dividido en dos capas:

\begin{itemize}
  \item `codigo/app/schemas/state.py`: modelos Pydantic estrictos para validar
        el estado, configuraciones, métricas, errores y artefactos;
  \item `codigo/app/graph/state.py`: `TypedDict` usado por LangGraph,
        validador de estado e inicializador del caso CWRU.
\end{itemize}

El inicializador `create\_initial\_cwru\_state` crea una ejecución preparada para
generar el manifiesto del dataset. El estado resultante no contiene señales ni
arrays, solo rutas, contexto experimental y campos reservados para artefactos
futuros.

\section{Contratos Pydantic implementados}

Se han añadido tres grupos de contratos:

\begin{itemize}
  \item `codigo/app/schemas/dataset.py`: manifiesto del dataset, canales y
        metadatos de fallo;
  \item `codigo/app/schemas/agent\_decisions.py`: decisiones estructuradas del
        supervisor, limpiador, estructurador, modelador, evaluador y redactor;
  \item `codigo/app/schemas/executor\_results.py`: resultados comunes y
        especializados de ejecutores deterministas.
\end{itemize}

Estos contratos consolidan la regla arquitectónica de que los agentes producen
configuraciones o decisiones validables, mientras que los ejecutores devuelven
rutas, métricas, errores y referencias a artefactos, pero no datos pesados.

\section{Artefactos previstos}

\begin{itemize}
  \item `codigo/data/interim/cwru\_bearing/manifest.csv`;
  \item `codigo/data/interim/cwru\_bearing/profile.json`;
  \item `codigo/data/processed/cwru\_bearing/clean\_signals/`;
  \item `codigo/data/tensors/cwru\_bearing/windows\_features.csv`;
  \item `codigo/data/tensors/cwru\_bearing/windows\_raw.npz`;
  \item `codigo/data/tensors/cwru\_bearing/splits.json`;
  \item `codigo/reports/cwru\_bearing/`;
  \item `codigo/models/cwru\_bearing/`.
\end{itemize}

\section{Adaptadores de Entrada}

Para reducir el acoplamiento al formato CWRU se ha introducido una capa ligera
de adaptadores en `codigo/app/services/signal\_adapters.py`. Esta capa expone
dos operaciones comunes: leer todos los canales numéricos de un fichero y cargar
un canal concreto como vector. Los ejecutores de perfilado y limpieza usan ahora
esta frontera en lugar de leer directamente ficheros `.mat`.

La implementación inicial soporta `.mat`, `.csv`, `.txt`, `.tsv` y `.npz`. En
el caso de CWRU, el adaptador reconoce los canales `DE\_time`, `FE\_time`,
`BA\_time` y `RPM`; en formatos tabulares, interpreta columnas numéricas como
canales. Para incorporar un nuevo origen, como NASA IMS, el objetivo será crear
o ampliar un adaptador de entrada y mantener estable el resto del pipeline:
perfilado, limpieza, estructuración, modelado y evaluación.

\section{Ejecutor de manifiesto}

Se ha implementado el primer ejecutor determinista del pipeline, el módulo
`dataset\_manifest.py`. Su responsabilidad es generar el artefacto:

\begin{center}
\begin{minipage}{0.85\textwidth}
\begin{verbatim}
codigo/data/interim/cwru_bearing/manifest.csv
\end{verbatim}
\end{minipage}
\end{center}

a partir de los ficheros `.mat` crudos de CWRU.

El ejecutor no carga señales completas. Solo inspecciona los nombres de fichero
y aplica un mapeo controlado de metadatos CWRU: carga, RPM, tipo de fallo,
diámetro y posición del fallo exterior cuando corresponde. El manifiesto
generado conserva la frecuencia original de cada fichero y fija una frecuencia
objetivo de 12 kHz para el MVP.

La ejecución real sobre los datos descargados ha generado 64 registros: 4
ficheros normales y 60 ficheros con fallo. El fichero resultante contiene 65
líneas incluyendo la cabecera.

\section{Ejecutor de perfilado}

Se ha implementado el perfilador determinista `data\_profiler.py`. Este ejecutor
lee el manifiesto, inspecciona los ficheros `.mat` y genera:

\begin{center}
\begin{minipage}{0.85\textwidth}
\begin{verbatim}
codigo/data/interim/cwru_bearing/profile.json
\end{verbatim}
\end{minipage}
\end{center}

El perfil contiene metadatos y estadísticas por canal: número de muestras,
valores no finitos, mínimo, máximo, media, desviación típica, RMS, curtosis y
energía. No almacena arrays ni señales completas.

La ejecución real ha perfilado 64 ficheros. El resumen agregado contiene 4
ficheros normales, 60 ficheros con fallo, frecuencias de muestreo de 48 kHz y
12 kHz, y los canales detectados `BA\_time`, `DE\_time`, `FE\_time` y `RPM`.

\section{Ejecutor de Limpieza}

Se ha implementado el ejecutor determinista `cleaning.py`. Este ejecutor lee el
manifiesto y el perfil, extrae el canal principal configurado, elimina valores
no finitos, remuestrea la señal a 12 kHz cuando procede y guarda una señal
limpia comprimida por fichero:

\begin{center}
\begin{minipage}{0.85\textwidth}
\begin{verbatim}
codigo/data/processed/cwru_bearing/clean_signals/
\end{verbatim}
\end{minipage}
\end{center}

También genera un resumen de auditoría en:

\begin{center}
\begin{minipage}{0.85\textwidth}
\begin{verbatim}
codigo/data/processed/cwru_bearing/cleaning_summary.json
\end{verbatim}
\end{minipage}
\end{center}

La ejecución real ha producido 64 ficheros `.npz` comprimidos. Los cuatro
ficheros normales se han remuestreado de 48 kHz a 12 kHz; los ficheros de fallo
ya estaban a 12 kHz. El estado del grafo conserva únicamente la ruta del
directorio limpio y las referencias a artefactos, no las señales.

\section{Ejecutor de Estructuración Temporal}

Se ha implementado el ejecutor determinista `structuring.py`. Este ejecutor lee
las señales limpias, valida que el canal y la frecuencia coincidan con la
configuración, genera ventanas temporales con tamaño 2048 y solapamiento del
50\%, calcula features temporales iniciales y crea particiones reproducibles a
nivel de fichero.

\begin{center}
\begin{minipage}{0.85\textwidth}
\begin{verbatim}
codigo/data/tensors/cwru_bearing/windows_features.csv
codigo/data/tensors/cwru_bearing/windows_raw.npz
codigo/data/tensors/cwru_bearing/splits.json
\end{verbatim}
\end{minipage}
\end{center}

La estrategia de partición evita mezclar ventanas de un mismo fichero entre
subconjuntos. Las ventanas normales se reparten entre entrenamiento,
validación y test, mientras que las ventanas con fallo se reservan para test.
La ejecución real ha generado 7460 ventanas de 2048 muestras: 175 de
entrenamiento, 117 de validación y 7168 de test. El tensor crudo tiene forma
`(7460, 2048)` y el directorio de tensores ocupa aproximadamente 22 MB.

\section{Ejecutor de Modelado Base}

Se ha implementado el ejecutor determinista `modeling.py` con el primer modelo
base de detección de anomalías: Isolation Forest. El modelo se entrena
únicamente con ventanas normales del subconjunto de entrenamiento y consume las
features temporales generadas por la fase anterior.

\begin{center}
\begin{minipage}{0.85\textwidth}
\begin{verbatim}
codigo/models/cwru_bearing/isolation_forest.joblib
codigo/models/cwru_bearing/predictions.csv
codigo/models/cwru_bearing/modeling_summary.json
\end{verbatim}
\end{minipage}
\end{center}

El artefacto `joblib` conserva el modelo entrenado, las columnas usadas y el
umbral de decisión. Las predicciones guardan, por cada ventana, identificador,
fichero original, partición, etiqueta real, puntuación de anomalía, umbral y
predicción binaria. Las métricas finales no se calculan en este ejecutor, sino
que se reservan para la fase de evaluación.

La ejecución real ha entrenado el modelo con 175 ventanas normales y ha
generado 7460 predicciones. El umbral estimado ha sido 0.6334 y el resumen de
modelado registra 403 ventanas predichas como normales y 7057 como anómalas.
Estos conteos son una salida preliminar del modelo, no una evaluación final de
rendimiento.

\section{Ejecutor de Evaluación}

Se ha implementado el ejecutor determinista `evaluation.py`. Este módulo lee las
predicciones del modelo, compara la etiqueta binaria real con la predicción de
anomalía y calcula métricas técnicas por partición. La partición principal para
el MVP es `test`.

\begin{center}
\begin{minipage}{0.85\textwidth}
\begin{verbatim}
codigo/reports/cwru_bearing/evaluation/metrics.json
codigo/reports/cwru_bearing/evaluation/evaluation_summary.md
\end{verbatim}
\end{minipage}
\end{center}

La ejecución real sobre el split de test ha obtenido precision 0.9997, recall
1.0000, F1-score 0.9999, ROC-AUC 1.0000, PR-AUC 1.0000 y una tasa de falsos
positivos de 0.0171. La matriz de confusión ha sido: 115 verdaderos negativos,
2 falsos positivos, 0 falsos negativos y 7051 verdaderos positivos.

Estos resultados deben interpretarse como validación del pipeline y del
baseline inicial sobre un benchmark controlado. No constituyen todavía una
garantía de generalización industrial, que se deberá analizar con validaciones
más exigentes y, posteriormente, con otros datasets como NASA IMS.

\section{Dependencias}

Se ha creado `requirements.txt` como punto inicial para construir el entorno de
Miniconda `tfm\_v2`, fijado sobre Python 3.11. Incluye dependencias de contratos
y grafo, procesamiento de datos, modelos clásicos, visualización técnica y
verificación.

\section{Verificación Actual}

La verificación de los contratos se ha realizado con tests unitarios en el
entorno objetivo:

\begin{center}
\begin{minipage}{0.85\textwidth}
\begin{verbatim}
conda run -n tfm_v2 python -m unittest discover codigo/tests
\end{verbatim}
\end{minipage}
\end{center}

Las pruebas validan la creación del estado inicial CWRU, serialización JSON,
rechazo de campos no declarados, restricciones de configuraciones y métricas,
coherencia del manifiesto, decisiones de agentes, resultados de ejecutores y
generación del manifiesto, perfilado, limpieza, estructuración temporal y
modelado base y evaluación CWRU. También validan los adaptadores de entrada para
formatos `.mat` y `.csv`. En esta fase se han ejecutado 48 tests
correctamente.

\section{Siguiente Paso de Implementación}

El siguiente avance técnico será montar el grafo LangGraph mínimo, encadenando
los ejecutores deterministas y preparando los puntos donde después entrarán los
agentes especializados.

\chapter{Experimentos}

\section{Protocolo experimental}

Esta sección documentará los experimentos ejecutados sobre CWRU Bearing Dataset.
Cada experimento deberá indicar:

\begin{itemize}
  \item versión del manifiesto;
  \item configuración de ventanas;
  \item features o tensores usados;
  \item modelo e hiperparámetros;
  \item estrategia de partición;
  \item métricas obtenidas;
  \item artefactos generados.
\end{itemize}

\section{Experimento base}

El primer experimento comparará modelos clásicos de detección de anomalías
entrenados con ventanas normales y evaluados contra ventanas normales retenidas
y ventanas con fallo.

\section{Experimentos posteriores}

Experimentos previstos:

\begin{itemize}
  \item sensibilidad al tamaño de ventana;
  \item comparación entre features temporales y frecuenciales;
  \item validación leave-one-load-out;
  \item validación sobre NASA IMS cuando el MVP esté estabilizado.
\end{itemize}

\chapter{Resultados}

\section{Resultados cuantitativos}

Esta seccion recogera las metricas de cada experimento.

\begin{table}[h]
\centering
\begin{tabular}{lrrrrr}
\toprule
Modelo & Precision & Recall & F1 & ROC-AUC & PR-AUC \\
\midrule
Pendiente & -- & -- & -- & -- & -- \\
\bottomrule
\end{tabular}
\caption{Plantilla de resultados para modelos de deteccion de anomalias.}
\label{tab:resultados_modelos}
\end{table}

\section{Analisis cualitativo}

Esta seccion interpretara falsos positivos, falsos negativos, sensibilidad a la
carga del motor, robustez por tipo de fallo y trazabilidad de artefactos.

\section{Limitaciones observadas}

Cada resultado debera acompanarse de limitaciones, supuestos y riesgos de
generalizacion industrial.

\chapter{Conclusiones}

\section{Conclusiones principales}

Esta sección resumirá las aportaciones del proyecto, separando los resultados
del pipeline, la arquitectura multiagente y los modelos de detección de
anomalías.

\section{Trabajo futuro}

Líneas previstas:

\begin{itemize}
  \item ampliar la validación con NASA IMS Bearing Dataset;
  \item incorporar autoencoders para ventanas crudas;
  \item añadir persistencia de checkpoints;
  \item exponer el MVP mediante API;
  \item preparar ejecución Docker y SLURM cuando el flujo local sea estable.
\end{itemize}


\begingroup
\sloppy
\bibliographystyle{plain}
\bibliography{bibliografia/referencias}
\endgroup

\end{document}
